\section{Conclusion}

This consultation constructed, optimized, and stress-tested a six-asset portfolio for Group 3 (Charlie)---TEL, MER, AEV, NVDA, META, and BTC---over the sample period January 2, 2020 through December 31, 2025, benchmarked against SPY at a 5.25 percent Philippine 91-day Treasury Bill risk-free rate.

The equal-weighted Buy-and-Hold strategy compounded USD 1.00 into USD 8.43 (743.04 percent cumulative, 42.83 percent CAGR) but ended as a concentrated two-asset bet, with NVDA and BTC reaching 61.75 percent and 24.77 percent of the portfolio and a -58.21 percent peak drawdown. Monthly rebalancing delivered USD 4.85 (384.85 percent, 30.21 percent CAGR) with materially lower volatility (24.18 percent versus 35.12 percent) and a shallower drawdown (-43.20 percent), demonstrating the discipline's risk-control value at the cost of some upside.

Mean-variance optimization formalized the trade-off. The GMV portfolio (TEL 31.01 percent, MER 31.23 percent, AEV 16.85 percent, NVDA 3.15 percent, META 12.86 percent, BTC 4.90 percent) achieves 16.03 percent return at 20.22 percent volatility (Sharpe 0.5332). The Maximum Sharpe portfolio (NVDA 55.44 percent, BTC 26.14 percent, MER 18.42 percent) achieves 58.10 percent return at 38.25 percent volatility and a Sharpe ratio of 1.3815, verified across multiple solver starts and cross-checked against a quadratic-programming frontier sweep.

Risk assessment showed that tail risk is real and non-Gaussian. The Maximum Sharpe portfolio's daily 99 percent Historical Expected Shortfall of 8.55 percent exceeds its Gaussian parametric estimate by 236 basis points, and its worst historical stress loss reaches -28.83 percent (COVID) with -45.30 percent under the 1.5x Black Swan scenario. The GMV portfolio generally lowers stress losses relative to the Maximum Sharpe portfolio, though the reduction is scenario-dependent.

The recommendation follows from these results: adopt the Maximum Sharpe portfolio under monthly rebalancing for investors with high loss capacity, or the GMV portfolio for capital preservation, with historical tail-risk measures governing the risk budget. All quantitative claims in this report are reproducible from the canonical R pipeline (QMD) and the frozen numerical summary, and every figure and table reflects the final analytical state.
